% ============================================================
% Chapter 6: Experimental Design
%   Section: Selection and Justification of the Evaluation Dataset
%     Subsection: The Closed-Retrieval Paradigm: FEVER
% Included from chapters/06-experimental-design/02-evaluation-dataset/evaluation-dataset.tex
% ============================================================

\subsection{The Closed-Retrieval Paradigm: FEVER}
\label{subsec:fever}

FEVER (\textit{Fact Extraction and VERification}) \cite{thorne2018fever}
is, to date, the most influential dataset in automated
fact-verification research. It consists of 185,445 claims generated
through the controlled alteration of sentences extracted from
Wikipedia, each subsequently labelled by human annotators ---
without knowledge of the original sentence --- as
\texttt{SUPPORTED}, \texttt{REFUTED}, or \texttt{NOT ENOUGH INFO}.
For each supported or refuted claim, the sentence or set of
sentences from Wikipedia constituting the necessary evidence for
that verdict is additionally recorded.

The aspect of FEVER most relevant to this work is its
\textbf{closed-domain retrieval paradigm}: both the claim and its
correct evidence originate from a fixed, bounded corpus --- a static
dump of English Wikipedia --- and systems evaluated against this
benchmark typically index that same corpus to perform evidence
retrieval. The original paper itself characterises the task as the
combination of three subtasks evaluated jointly: document selection,
sentence selection within those documents, and final verdict
classification, all operating over the same closed corpus.

Direct multilingual extensions of FEVER also exist, such as XFEVER,
which automatically translate (via DeepL) the original claim-evidence
pairs into other languages, including Spanish. However, these
extensions inherit the same structural limitation as the original
dataset: the \textit{correct} evidence remains defined as a specific
sentence within a closed, translated corpus, rather than as any
valid evidence available on the open web.
